% Part: turing-machines
% Chapter: undecidability
% Section: enumerating-tms

\documentclass[../../../include/open-logic-section]{subfiles}

\begin{document}

\olfileid{tur}{und}{enu}
\olsection{برشمردن ماشین‌های تورینگ}

\begin{explain}
می‌توانیم نشان دهیم مجموعهٔ همهٔ ماشین‌های تورینگ !!{enumerable} است.
این از آنجا نتیجه می‌شود که هر ماشین تورینگ را می‌توان به‌طور متناهی
توصیف کرد. مجموعهٔ حالت‌ها و واژگان نوار مجموعه‌هایی متناهی‌اند. تابع
گذار تابعی جزئی از~$Q\times\Sigma$ به
$Q\times\Sigma\times\{\TMleft,\TMright,\TMstay\}$ است و آن را نیز
می‌توان با فهرست‌کردن مقادیرش برای تعداد متناهی زوج‌آرگومانی که روی
آن‌ها تعریف شده است مشخص کرد.

این سخن تا همین‌جا درست است، اما تفاوت ظریفی وجود دارد. تعریف
ماشین‌های تورینگ هیچ محدودیتی بر !!{element}sیی که مجموعهٔ حالت‌ها و
الفبای نوار می‌توانند داشته باشند نگذاشت. پس، برای نمونه، از نظر فنی
به ازای هر عدد حقیقی ماشینی تورینگ وجود دارد که آن عدد را به‌عنوان
حالت به کار می‌برد. بااین‌حال \emph{رفتار} ماشین تورینگ مستقل از این
است که چه اشیائی نقش حالت‌ها و واژگان را داشته باشند. دو ماشین
تورینگِ \olref{fig:variants} را در نظر بگیرید.
\begin{figure}
\begin{center}
\begin{tikzpicture}[->,>=stealth',shorten >=1pt,auto,node distance=2.8cm,
                    semithick]
  \tikzstyle{every state}=[fill=none,draw=black,text=black]

  \node[initial,state]         (A)              {$q_0$};
  \node[state]         (B) [right of=A] {$q_1$};

  \path (A) edge [bend left] node {\TMtrans{\TMstroke}{\TMstroke}{\TMright}} (B)
        (B) edge [loop above] node {\TMtrans{\TMblank}{\TMblank}{\TMright}} (B)
            edge [bend left] node {\TMtrans{\TMstroke}{\TMstroke}{\TMright}} (A);
\end{tikzpicture}\\
\begin{tikzpicture}[->,>=stealth',shorten >=1pt,auto,node distance=2.8cm,
  semithick]
\tikzstyle{every state}=[fill=none,draw=black,text=black]

\node[initial,state]         (A)              {$s$};
\node[state]         (B) [right of=A] {$h$};

\path (A) edge [bend left] node {\TMtrans{A}{A}{\TMright}} (B)
(B) edge [loop above] node {\TMtrans{\TMblank}{\TMblank}{\TMright}} (B)
edge [bend left] node {\TMtrans{A}{A}{\TMright}} (A);
\end{tikzpicture}
\end{center}
\caption{گونه‌های ماشین \emph{زوج}}
\ollabel{fig:variants}
\end{figure}
این دو نمودار با دو ماشین متناظرند: $M$ با الفبای نوار
$\Sigma=\{\TMendtape,\TMblank,\TMstroke\}$ و مجموعه‌حالت‌های
$\{q_0,q_1\}$؛ و~$M'$ با الفبای
$\Sigma'=\{\TMendtape,\TMblank,A\}$ و حالت‌های~$\{s,h\}$. اما
دستورهایشان از هر جهت دیگر یکسان است: $M$ روی دنباله‌ای از~$n$ نماد
$\TMstroke$ متوقف می‌شود اگر و تنها اگر $n$ زوج باشد، و~$M'$ روی
دنباله‌ای از~$n$ نماد~$A$ متوقف می‌شود اگر و تنها اگر $n$ زوج باشد.
تنها کاری که کرده‌ایم تغییر نام~$\TMstroke$ به~$A$، $q_0$ به~$s$ و
$q_1$ به~$h$ است. این مثال تعمیم می‌یابد: می‌توانیم ماشین‌های تورینگ
را تا هنگامی یکسان بدانیم که یکی از دیگری با چنین تغییر نامی در
نمادها و حالت‌ها حاصل شود. در واقع، می‌توانیم نمادها و حالت‌های یک
ماشین تورینگ را صرفاً اعداد صحیح مثبت بدانیم: به‌جای~$\sigma_0$ عدد
$1$، به‌جای~$\sigma_1$ عدد~$2$ و به همین ترتیب؛ $\TMendtape$ عدد~$1$،
$\TMblank$ عدد~$2$ و مانند آن. بدین شیوه ماشین \emph{زوج} به ماشین
نشان‌داده‌شده در \olref{fig:standard-even} تبدیل می‌شود.
\begin{figure}
\[\begin{tikzpicture}[->,>=stealth',shorten >=1pt,auto,node distance=2.8cm,
  semithick]
\tikzstyle{every state}=[fill=none,draw=black,text=black]

\node[initial,state]         (A)              {$1$};
\node[state]         (B) [right of=A] {$2$};

\path (A) edge [bend left] node {\TMtrans{3}{3}{\TMright}} (B)
(B) edge [loop above] node {\TMtrans{2}{2}{\TMright}} (B)
edge [bend left] node {\TMtrans{3}{3}{\TMright}} (A);
\end{tikzpicture}
\]
\caption{یک ماشین \emph{زوج} استاندارد}
\ollabel{fig:standard-even}
\end{figure}
می‌توانیم ماشین تورینگی را که حالت‌ها و نمادهایش اعداد صحیح مثبت‌اند
ماشین \emph{استاندارد} بنامیم و ازاین‌پس فقط ماشین‌های استاندارد را
در نظر بگیریم.\footnote{اصطلاح «ماشین استاندارد» استاندارد نیست.}

می‌خواستیم نشان دهیم مجموعهٔ ماشین‌های تورینگ !!{enumerable} است و با
توجه به ملاحظات بالا کافی است نشان دهیم مجموعهٔ ماشین‌های تورینگ
استاندارد !!{enumerable} است. فرض کنید ماشین تورینگ استاندارد
$M=\tuple{Q,\Sigma,q_0,\delta}$ داده شده باشد. چگونه می‌توانیم آن را
با رشته‌ای متناهی از اعداد صحیح مثبت توصیف کنیم؟ نخست تعداد حالت‌ها،
خود حالت‌ها، تعداد نمادها، خود نمادها و حالت آغازین را فهرست می‌کنیم.
(به یاد آورید همهٔ این‌ها اعداد صحیح مثبت‌اند، زیرا $M$ ماشینی
استاندارد است.) دربارهٔ~$\delta$ چه می‌کنیم؟ مجموعهٔ آرگومان‌های ممکن،
یعنی زوج‌های~$\tuple{q,\sigma}$، متناهی است زیرا $Q$ و~$\Sigma$
متناهی‌اند. پس اطلاعات موجود در~$\delta$ صرفاً فهرست متناهی همهٔ
پنج‌تایی‌های~$\tuple{q,\sigma,q',\sigma',d}$ است که در آن
$\delta(q,\sigma)=\tuple{q',\sigma',D}$ و~$d$ عددی است که جهت~$D$ را
رمزگذاری می‌کند (مثلاً $1$ برای~$\TMleft$، $2$ برای~$\TMright$ و~$3$
برای~$\TMstay$).

بدین شیوه هر ماشین تورینگ استاندارد را می‌توان با فهرستی متناهی از
اعداد صحیح مثبت، یعنی دنباله‌ای
$s_M\in(\PosInt)^*$، توصیف کرد. برای نمونه، ماشین \emph{زوج}
استاندارد با دنبالهٔ زیر رمزگذاری می‌شود:
\[
2, \underbrace{1, 2}_Q, 3, \overbrace{1, 2, 3}^\Sigma, 1, \underbrace{1, 3, 2, 3, 2}_{\delta(1,3) = \tuple{2,3,R}},
\overbrace{2, 2, 2, 2, 2}^{\delta(2,2) = \tuple{2,2,R}},
\underbrace{2, 3, 1, 3, 2}_{\delta(2,3) = \tuple{1,3,R}}.
\]
\end{explain}

\begin{thm}
توابعی از~$\Nat$ به~$\Nat$ وجود دارند که تورینگ‌محاسبه‌پذیر نیستند.
\end{thm}

\begin{proof}
می‌دانیم مجموعهٔ دنباله‌های متناهیِ اعداد صحیح مثبت~$(\PosInt)^*$
!!{enumerable} است (\cref{sfr:siz:zigzag:prob:posint-star}). ازاین‌رو
مجموعهٔ توصیف‌های ماشین‌های تورینگ استاندارد نیز، به‌عنوان زیرمجموعه‌ای
از~$(\PosInt)^*$، برشمردنی است. هر تابع تورینگ‌محاسبه‌پذیر از~$\Nat$
به~$\Nat$ را یک ماشین تورینگ (در واقع ماشین‌های فراوانی) محاسبه
می‌کند. با تغییر نام حالت‌ها و نمادهای آن به اعداد صحیح مثبت (به‌ویژه
$\TMendtape$ به~$1$، $\TMblank$ به~$2$ و~$\TMstroke$ به~$3$) می‌بینیم
که هر تابع تورینگ‌محاسبه‌پذیر را یک ماشین استاندارد محاسبه می‌کند.
پس مجموعهٔ همهٔ توابع تورینگ‌محاسبه‌پذیر از~$\Nat$ به~$\Nat$ نیز
برشمردنی است.

از سوی دیگر، مجموعهٔ همهٔ توابع از~$\Nat$ به~$\Nat$ !!{enumerable}
نیست (\cref{sfr:siz:red:prob:nat-nat}). اگر همهٔ توابع را ماشینی
تورینگ محاسبه می‌کرد، می‌توانستیم با فهرست‌کردن همهٔ توصیف‌های
ماشین‌های تورینگی که آن‌ها را محاسبه می‌کنند مجموعهٔ همهٔ توابع را
برشماریم. بنابراین توابعی وجود دارند که تورینگ‌محاسبه‌پذیر نیستند.
\end{proof}

\begin{prob}
  آیا راهی برای توصیف ماشین‌های تورینگ می‌یابید که نیازی به
  فهرست‌کردن صریح حالت‌ها و نمادهای الفبا نداشته باشد؟ می‌توانید مفهوم
  خود را از ماشین «استاندارد» تعریف کنید، اما توضیح دهید چرا هر ماشین
  تورینگ را ماشینی «استاندارد» در معنای جدید شما می‌تواند شبیه‌سازی
  کند.
\end{prob}

\end{document}
